```latex
\documentclass[11pt]{amsart}
\usepackage{amsmath,amssymb,amsthm}
\usepackage[margin=1.15in]{geometry}

\newcommand{\M}{\mathcal{M}}
\newcommand{\T}{\mathcal{T}}
\newcommand{\C}{\mathbb{C}}
\newcommand{\Prym}{\mathrm{P}}
\DeclareMathOperator{\Hilb}{Hilb}
\DeclareMathOperator{\codim}{codim}

\theoremstyle{plain}
\newtheorem{theorem}{Theorem}[section]
\newtheorem{proposition}[theorem]{Proposition}
\newtheorem{lemma}[theorem]{Lemma}
\newtheorem{corollary}[theorem]{Corollary}
\theoremstyle{definition}
\newtheorem{definition}[theorem]{Definition}
\newtheorem{claim}[theorem]{Claim}
\theoremstyle{remark}
\newtheorem{remark}[theorem]{Remark}

\begin{document}

\title{On an asserted boundary rigidity for deformations of\\ totally geodesic surfaces in the moduli space of curves}
\author{Anonymous}
\date{\today}

\begin{abstract}
We analyze the following research problem. Let $V_0\subset \M_g$ be a smooth complex
surface which is totally geodesic for the Teichm\"uller metric, let $(V_t)_{t\in\Delta}$
be a holomorphic family deforming it, subject to certain regularity hypotheses on the
intersections of the closures $\overline{V_t}\subset\overline{\M}_g$ with the
Deligne--Mumford boundary strata, and suppose that for some boundary stratum
$\Delta_\Gamma\cong \M_{g_1,n_1}\times\cdots\times \M_{g_k,n_k}$ the trace
$\partial V_t=\overline{V_t}\cap\Delta_\Gamma$ is a curve which is collapsed by one of
the projections $\pi_i$ to a single point $Y_t$. The problem asks for a proof that
$Y_t=Y_0$ for all $t$ (``non-varying'').

We prove everything in this setup that admits proof: the assignment $t\mapsto Y_t$ is
holomorphic (Lemma~\ref{lem:hol}); all hypotheses other than the total geodesy of the
central fiber are compatible with a non-constant $Y_t$ (Theorem~\ref{thm:model});
under mild genericity hypotheses, first-order deformations with non-zero infinitesimal
motion of the collapse point always exist (Theorem~\ref{thm:inflect}); and in the
algebraic category the traces must be complete curves, which renders the hypotheses
contradictory in a number of low-complexity cases (Proposition~\ref{prop:complete}).
We then show (Theorem~\ref{thm:dichotomy}) that the asserted statement is subject to
the following dichotomy: either it is vacuous---because no totally geodesic complex
surface in any $\M_g$ has ever been constructed and their existence appears to be an
open problem---or it asserts an extrinsic deformation rigidity for such surfaces
(Proposition~\ref{prop:R}) for which no mechanism exists in the stated hypotheses or
anywhere in the literature. In neither branch can a complete rigorous proof be given
at present. We therefore report, with full supporting analysis, that the requested
proof does not exist, and we explain precisely why any purported proof must fail.
\end{abstract}

\maketitle

\section{Introduction and statement of the problem}\label{sec:intro}

\subsection*{The claim under discussion}

Let $\M_g$ denote the coarse moduli space of smooth complex curves of genus $g$
($g\ge 2$), and $\overline{\M}_g$ its Deligne--Mumford compactification
\cite{DM}, whose boundary decomposes into strata $\Delta_G$ indexed by stable dual
graphs $G$. The problem, as posed to us, is the following.

\begin{claim}[Claim $\mathsf{P}$]\label{claim:P}
Let $V_0\subset\M_g$ be a smooth complex subvariety of dimension $2$ which is totally
geodesic with respect to the Teichm\"uller metric. Let $(V_t)_{t\in\Delta}$ be a
holomorphic family of submanifolds deforming $V_0$ (in the sense of the problem's
footnote: a complex submanifold $\mathcal V\subset \M_g\times\Delta$ with
$V_t=\psi(\phi^{-1}(t))$ for the natural projections, the family being smoothly
trivial), satisfying:
\begin{enumerate}
\item[(i)] for every boundary stratum $\Delta_G$, the family
$\overline{V_t}\cap \Delta_G$ is a (smoothly trivial) family; and
\item[(ii)] the induced incidence pattern with pairs of strata
$\Delta_G \supset \Delta_{G'}$ agrees, for all $t$, with the pattern at $t=0$.
\end{enumerate}
Note that $V_t$ is \emph{not} assumed totally geodesic for $t\neq 0$.
Fix a stable graph $\Gamma$ such that
$\Delta_\Gamma\cong \M_{g_1,n_1}\times\cdots\times \M_{g_k,n_k}$ (finite group actions
excluded), with projections $\pi_i$, and write $\partial V_t:=\overline{V_t}\cap
\Delta_\Gamma$. Suppose that $\dim \partial V_t=1$ and that for some fixed index $i$
one has $\pi_i(\partial V_t)=\{Y_t\}$, a single point, for every $t$.
\emph{Conclusion to be proved:} $Y_t=Y_0$ for all $t\in\Delta$; i.e., the family
$\partial V_t$ is ``non-varying'' on the factor $\M_{g_i,n_i}$.
\end{claim}

We were instructed to work hard and to return only with a complete rigorous proof of
Claim~\ref{claim:P}. After a careful analysis we must report that \emph{no such proof
exists in the present state of mathematical knowledge}, and that a purported proof
would necessarily be erroneous. The purpose of this paper is to substantiate this
report with complete proofs of all the assertions needed to make it precise. Our
findings are:

\begin{itemize}
\item The conclusion of Claim~\ref{claim:P} is equivalent to the constancy of a
canonically associated holomorphic map $Y:\Delta\to \M_{g_i,n_i}$
(Lemma~\ref{lem:hol}); it is therefore a genuine rigidity statement, not a formal
consequence of the definitions.
\item All the hypotheses of Claim~\ref{claim:P} \emph{except} the total geodesy of
$V_0$ are satisfied, in an explicitly constructed model, by families for which
$Y_t$ moves non-trivially (Theorem~\ref{thm:model}). Hence no argument based on the
regularity hypotheses (i), (ii), on smooth triviality, or on the dimension and
collapse conditions, can prove Claim~\ref{claim:P}: the total geodesy of the single
central fiber must do all of the work.
\item Under mild, explicit genericity hypotheses, first-order deformations with
non-zero infinitesimal velocity of the collapse point always exist
(Theorem~\ref{thm:inflect}). Thus Claim~\ref{claim:P}, if true for a given totally
geodesic surface, amounts to the assertion that none of these first-order
deformations integrates within the class of collapsed traces
(Proposition~\ref{prop:R}): a rigidity of a type for which no mechanism is known.
\item The objects to which Claim~\ref{claim:P} applies are not known to exist: to the
best of our knowledge, \emph{no totally geodesic complex surface in any $\M_g$ has
ever been constructed} (Section~\ref{sec:knowledge}). The existence question is open.
Consequently Claim~\ref{claim:P} is either vacuously true (with vacuity equivalent to
an open problem) or it asserts the unprecedented rigidity of
Proposition~\ref{prop:R} for hypothetical objects.
\end{itemize}

\subsection*{Standing interpretations}\label{subsec:standing}

As worded, Claim~\ref{claim:P} leaves several analytic structures unspecified, and
one must fix a reading before the problem is well posed. We adopt the following
minimal reading; our negative conclusions hold \emph{a fortiori} under weaker
readings.

\begin{definition}[Standing Interpretation]\label{def:standing}
Throughout, $\Delta$ is the open unit disc, and we assume:
\begin{enumerate}
\item[$(\alpha)$] $\mathcal V\subset \M_g\times\Delta$ is a complex submanifold,
$\phi:\mathcal V\to\Delta$ is a holomorphic submersion, and the closure
$\overline{\mathcal V}$ of $\mathcal V$ in $\overline{\M}_g\times\Delta$ is a closed
analytic subset whose fiber over $t$ equals $\overline{V_t}$.
\item[$(\beta)$] For each stable graph $G$, the set
$\mathcal B_G:=\overline{\mathcal V}\cap(\Delta_G\times\Delta)$, with its reduced
analytic structure, is a complex space, closed in $\Delta_G\times\Delta$, and the
projection $\mathcal B_G\to\Delta$ is a proper holomorphic submersion. This is the
precise meaning we assign to the phrase ``(smoothly trivial) family'' in hypothesis (i).
\end{enumerate}
If hypothesis (i) is read more weakly (mere $C^\infty$ triviality of the trace
families), then even the holomorphicity of $t\mapsto Y_t$---a prerequisite for the
claim to have definite content---is not available without additional assumptions,
and the problems identified below are only aggravated.
\end{definition}

\begin{remark}[Stack issues]
To avoid notational overhead we ignore the standard stack-theoretic subtleties of
$\overline{\M}_g$: at stable curves without automorphisms the stack is locally
isomorphic to its coarse space, and all our local arguments are performed at such
points. Global arguments may be performed after passage to a smooth finite cover
(e.g.\ a level cover), which changes nothing in the analysis. See
\cite{DM, Knudsen}.
\end{remark}

\subsection*{Structure of the paper}

Section~\ref{sec:struct} records elementary structural facts.
Section~\ref{sec:hol} proves the holomorphicity of the collapse map.
Section~\ref{sec:model} constructs the model family showing that the regularity
hypotheses are compatible with a moving collapse point, and derives the constraints
on traces in the algebraic category. Section~\ref{sec:inflect} performs the
first-order analysis. Section~\ref{sec:knowledge} reviews the state of knowledge
about totally geodesic subvarieties and states the dichotomy.
Section~\ref{sec:conc} states our conclusion.

\section{Structural preliminaries}\label{sec:struct}

\begin{lemma}\label{lem:divisor}
Assume $\overline{V_0}$ is irreducible of dimension $2$ and meets $\Delta_\Gamma$
non-emptily and properly. Then $\Delta_\Gamma$ is a boundary divisor; in the
situation of Claim~\ref{claim:P} (where $\Delta_\Gamma$ is, up to a finite group, a
product of moduli spaces of pointed curves), $\Gamma$ has a single separating edge
and $\Delta_\Gamma\cong \M_{g_1,1}\times\M_{g_2,1}$ with $g_1+g_2=g$.
\end{lemma}

\begin{proof}
For irreducible subvarieties $A,B$ of a smooth variety with $A\cap B\neq\emptyset$,
one has $\dim(A\cap B)\ge \dim A+\dim B-\dim X$ always, with equality for proper
intersections. Writing $N=3g-3$ and $c=\operatorname{codim}\Delta_\Gamma$, properness
gives $1=\dim(\overline{V_0}\cap\Delta_\Gamma)=2+(N-c)-N=2-c$, so $c=1$. A
boundary stratum of codimension $1$ corresponds to a stable graph with one edge; the
product structure $\prod_j \M_{g_j,n_j}$ with $k\ge 2$ factors forces the edge to be
separating, and then the stratum is (up to the finite group excluded by hypothesis)
$\M_{g_1,1}\times\M_{g_2,1}$, the identification being the canonical one (normalize
the one-nodal curve at its node).
\end{proof}

\begin{remark}\label{rem:deep}
If instead the intersection is non-proper, then $\Gamma$ may have several edges and
$\overline{V_0}$ lies in special position with respect to $\Delta_\Gamma$. The
hypotheses of Claim~\ref{claim:P} allow this; it only makes the asserted conclusion
less plausible, since the configuration is then less constrained, and none of our
arguments below require properness.
\end{remark}

\section{The collapse point varies holomorphically}\label{sec:hol}

\begin{lemma}\label{lem:hol}
Under the Standing Interpretation~\ref{def:standing} and the hypotheses of
Claim~\ref{claim:P}, the assignment $t\mapsto Y_t$ is a holomorphic map
$Y:\Delta\to \M_{g_i,n_i}$.
\end{lemma}

\begin{proof}
Consider $\mathcal B_\Gamma=\overline{\mathcal V}\cap(\Delta_\Gamma\times\Delta)$
with its reduced analytic structure; by $(\beta)$ the projection
$p:\mathcal B_\Gamma\to\Delta$ is a proper holomorphic submersion, and by $(\alpha)$
its fiber over $t$ is $\partial V_t$, a compact complex curve. Let
$q:\mathcal B_\Gamma\to\Delta_\Gamma$ be the second projection and put
$\widetilde Y:=\pi_i\circ q$, a holomorphic map. By the collapse hypothesis, for
every $t$ and every $b\in\mathcal B_{\Gamma,t}=\partial V_t$ we have
$\widetilde Y(b)\in \pi_i(\partial V_t)=\{Y_t\}$; hence $\widetilde Y$ is constant on
each fiber of $p$.

By the Fischer--Grauert theorem \cite{FG}, a proper holomorphic submersion is locally
trivial: for each $t_0\in\Delta$ there exist a neighborhood $U\ni t_0$ and a
biholomorphism $\mathcal B_\Gamma|_U\cong U\times F$ over $U$. In this trivialization
$\widetilde Y(u,f)$ is independent of $f$, so the formula
$Y(u):=\widetilde Y(u,f_0)$ (for any fixed $f_0\in F$) defines a holomorphic function
of $u$ which is independent of the choices, because fibers of $p$ map to single
points. Patching over $\Delta$ gives a holomorphic $Y:\Delta\to\M_{g_i,n_i}$ with
$\pi_i(\partial V_t)=\{Y(t)\}$.
\end{proof}

\begin{corollary}\label{cor:equiv}
Claim~\ref{claim:P} is equivalent, for each admissible family, to the assertion that
the holomorphic map $Y$ of Lemma~\textup{\ref{lem:hol}} is constant.
\end{corollary}

\begin{proof}
Since $\Delta$ is connected and $\M_{g_i,n_i}$ is a (reduced, Hausdorff) complex
space, a holomorphic map from $\Delta$ is constant if and only if it is locally
constant. The conclusion of Claim~\ref{claim:P} is precisely $Y\equiv Y(0)$.
\end{proof}

Thus Claim~\ref{claim:P} is a rigidity statement about the holomorphic map $Y$, and
any proof must exhibit a mechanism forcing $Y$ to be locally constant.

\section{The auxiliary hypotheses are compatible with a moving collapse point}
\label{sec:model}

The goal of this section is to show that the hypotheses (i), (ii), the smooth
triviality of the family, the dimension condition $\dim\partial V_t=1$, and the
collapse condition $\pi_i(\partial V_t)=\{Y_t\}$, \emph{jointly} place no constraint
whatsoever on $Y$ beyond its holomorphicity: there are configurations satisfying all
of them in which $Y$ is non-constant. Consequently, any proof of
Claim~\ref{claim:P} would have to derive the constancy of $Y$ from the total geodesy
of $V_0$ alone.

\subsection{Local coordinates near a separating node}

\begin{lemma}\label{lem:coords}
Let $g=3$, let $\Gamma$ be the stable graph with one separating edge joining genera
$1$ and $2$, so that $\Delta_\Gamma\cong \M_{1,1}\times\M_{2,1}$ with projections
$\pi_1,\pi_2$. Let $X_0=E_0\cup C_0$ be a stable curve of type $\Gamma$ whose two
pointed components $(E_0,p_0)$ and $(C_0,q_0)$ are smooth and automorphism-free, and
let $x_0=[X_0]\in\Delta_\Gamma$. Then there exist an open neighborhood $W$ of $x_0$
in $\overline{\M}_3$ and a holomorphic coordinate chart
\[
\chi=(m_1,m_2,u):\quad W\longrightarrow \C\oplus\C^4\oplus\C ,
\]
with $\chi(x_0)=0$, such that:
\begin{enumerate}
\item $\chi(W\cap\Delta_\Gamma)=\{u=0\}$;
\item $W$ meets no boundary stratum other than $\Delta_\Gamma$ (in particular
$W\cap(\Delta_{G})=\emptyset$ for all $G\neq\Gamma$);
\item on $\{u=0\}$, $\pi_1\circ\chi^{-1}$ is the projection $(m_1,m_2,0)\mapsto m_1$,
where $m_1$ is a local coordinate on $\M_{1,1}$;
\item every point of $W$ with $u\neq 0$ lies in $\M_3$.
\end{enumerate}
\end{lemma}

\begin{proof}
Since $X_0$ has a single node and smooth components, $x_0$ lies in the smooth locus
of the coarse space $\overline{\M}_3$, and the boundary is locally a smooth divisor
\cite{DM,Knudsen}. By the local structure theory of the moduli of stable curves
(the deformation space of a compact-type nodal curve factors canonically as the
product of the deformation spaces of the two pointed components and the one-dimensional
smoothing space of the node; see \cite[Ch.~XII]{ACGII} and \cite{Sernesi}), and by the
canonical identification $\Delta_\Gamma\cong \M_{1,1}\times\M_{2,1}$ (normalization
at the node), there are local coordinates $(m_1,m_2,u)$ as claimed in which the
boundary is $\{u=0\}$ and the $m$-coordinates restrict on the boundary to coordinates
on the two factors. Assertion (2) holds because every other boundary stratum is a
closed subset of $\overline{\M}_3$ not containing $x_0$ (a curve with one node is
not in the closure of a stratum whose generic point has more nodes, nor in the
closure of the other one-edge strata), so their complement is an open neighborhood of
$x_0$. Assertion (4) follows from (1)--(2): a point of $W$ with $u\ne0$ cannot lie in
$\Delta_\Gamma$ by (1), nor in any other stratum by (2), hence lies in $\M_3$.
\end{proof}

\subsection{The model family}

\begin{theorem}\label{thm:model}
In the setting of Lemma~\textup{\ref{lem:coords}}, let
$v\in\C\setminus\{0\}$ and $w\in\C^4\setminus\{0\}$, and let
$\rho_u>0$ be small enough that the relevant closed polydiscs lie in the chart
domain. For $t\in\Delta_{1/2}:=\{|t|<\tfrac12\}$ define
\[
V_t:=\chi^{-1}\Big(\big\{(t v,\; s w,\; u)\ :\ s\in\Delta_{1/2},\ 0<|u|<\rho_u\big\}\Big)
\subset \M_3 .
\]
Then:
\begin{enumerate}
\item $\mathcal V:=\{(x,t)\in \M_3\times\Delta_{1/2}\;:\; x\in V_t\}$ is a complex
submanifold of $\M_3\times\Delta_{1/2}$, and the projections realize a smoothly
trivial (indeed globally product-trivial) holomorphic family of two-dimensional
submanifolds deforming $V_0$ in the sense of the problem's footnote;
\item each $\overline{V_t}$ (closure in $\overline{\M}_3$) is compact, contained in
$W$, and meets the boundary \emph{only} in $\Delta_\Gamma$; hence hypothesis (i) of
Claim~\ref{claim:P} holds for all $G\neq\Gamma$ vacuously, and hypothesis (ii) holds
vacuously;
\item $\partial V_t=\overline{V_t}\cap\Delta_\Gamma$ is a compact one-dimensional
manifold with boundary, varying in a $C^\infty$-trivial (indeed fiberwise
holomorphic) family over $t$, so $\dim\partial V_t=1$;
\item $\pi_1(\partial V_t)=\{Y_t\}$, where $Y_t$ is the point of $\M_{1,1}$ with
local coordinate $t v$; the map $t\mapsto Y_t$ is holomorphic and non-constant.
\end{enumerate}
In particular, \emph{every} hypothesis of Claim~\ref{claim:P} is satisfied by this
family except the total geodesy of the central fiber $V_0$, while the conclusion of
Claim~\ref{claim:P} fails ($Y_t\ne Y_0$ for $t\neq 0$).
\end{theorem}

\begin{proof}
(1): $\mathcal V$ is the image of the holomorphic embedding
$\Theta:\Delta_{1/2,t}\times\Delta_{1/2,s}\times\Delta^{*}_{\rho_u}\to
\M_3\times\Delta_{1/2,t}$, $(t,s,u)\mapsto(\chi^{-1}(tv,sw,u),\,t)$. This map is
injective ($t$ is the second coordinate and $(s,u)$ is read off from $\chi$) and is
an immersion ($\chi$ is a chart), and it is a homeomorphism onto its locally closed
image; hence $\mathcal V$ is an embedded complex submanifold and $\phi$ is the
submersion $(x,t)\mapsto t$. The $(s,u)$-coordinates identify $\mathcal V$
globally with $\Delta_{1/2}\times(\Delta_{1/2,s}\times\Delta^*_{\rho_u})$, so the
family is product-trivial.

(2): The set $\chi^{-1}(\{tv\}\times\bar\Delta_{1/2,s}\times\bar\Delta_{\rho_u})$ is
compact, contained in $W$ for $\rho_u$ small, and contains all limit points of
sequences from $V_t$; hence it is $\overline{V_t}$. By Lemma~\ref{lem:coords}(1)--(2),
$\overline{V_t}\cap\Delta_G=\emptyset$ for $G\neq\Gamma$, and
$\overline{V_t}\cap\Delta_\Gamma=\chi^{-1}(\{tv\}\times\bar\Delta_{1/2,s}\times\{0\})$.

(3): This last set is a closed disc, a compact complex one-dimensional manifold with
boundary; the family of traces is the image of the holomorphic (in $(t,s)$) map
$(t,s)\mapsto \chi^{-1}(tv,sw,0)$, in particular $C^\infty$-trivial.

(4): On $\{u=0\}$ the projection $\pi_1$ reads off the $m_1$-coordinate
(Lemma~\ref{lem:coords}(3)), which is constantly $tv$ along the trace; hence
$\pi_1(\partial V_t)$ is the single point $Y_t$ of $\M_{1,1}$ with coordinate $tv$.
Since coordinates are injective, $Y_t=Y_0$ iff $t=0$, and $t\mapsto Y_t$ is
holomorphic and non-constant.
\end{proof}

\begin{remark}[Scope of Theorem~\ref{thm:model}]\label{rem:scope}
Two caveats delimit the category in which the model lives. First, the fibers $V_t$
are \emph{locally closed} but not closed submanifolds of $\M_3$ (their closures add
points at $|u|=\rho_u$, which lie in $\M_3$); this is consistent with the problem's
own footnote, which requires only a complex submanifold
$\mathcal V\subset \M_g\times\Delta$ with smoothly trivial projections. Second, the
trace families are compact manifolds with boundary and $C^\infty$-trivial families in
the sense of the parenthetical ``(smoothly trivial)'' in hypothesis (i); they are
not complex-analytic subsets of $\Delta_\Gamma$ (a closed disc is not analytic).
Both caveats concern only the strength of the category, not the logical point being
made, which is this: \emph{the totality of the auxiliary hypotheses of
Claim~\ref{claim:P} is invariant under, and compatible with, a non-constant
holomorphic collapse map $Y$.} Hence no argument can prove Claim~\ref{claim:P} by
operating on these hypotheses alone; the total geodesy of $V_0$ must be invoked, and
it must be invoked in a way that constrains \emph{all} deformations of $V_0$ as a
submanifold, not merely the constant family.
\end{remark}

\begin{remark}[The algebraic category]\label{rem:alg}
If one strengthens the reading so that the $V_t$ must be closed algebraic surfaces in
$\M_g$, then the traces acquire a rigid structure; this is analyzed in
Proposition~\ref{prop:complete} below. Constructing fully compliant examples in
that stronger category amounts to plumbing a constant family of pointed curves into
a complete family of pointed curves over a complete base (using Kodaira-type
families \cite{Kodaira} to produce complete curves in $\M_{g',1}$ for $g'\ge 3$,
together with the standard algebraic/analytic plumbing construction
\cite{Hubbard76, ACGII}). This construction is routine in outline but lengthy; since
none of our conclusions depends on it, we do not spell it out.
\end{remark}

\subsection{Constraints in the closed algebraic category}

\begin{proposition}\label{prop:complete}
Let $g\ge 3$, let $V\subset \M_g$ be a closed analytic subset of dimension $2$ such
that $\overline{V}$ is analytic in $\overline{\M}_g$ (automatic for closed analytic
$V$ by the Remmert--Stein extension theorem, since $\dim V=2<3g-4=\dim\partial\M_g$;
see \cite{Chirka, Demailly}), and let $\Gamma$ be such that
$\partial V:=\overline{V}\cap\Delta_\Gamma$ has pure dimension $1$. Then:
\begin{enumerate}
\item $\overline{V}$ is a projective algebraic surface, by Chow's theorem
\cite{Chirka, Demailly}.
\item Either $\partial V$ is compact---in which case it is a \emph{complete}
algebraic curve in $\Delta_\Gamma$---or else $\overline{V}$ meets deeper boundary
strata in the limit points of $\partial V$.
\item In the compact case, if moreover $\pi_i(\partial V)=\{Y\}$ for a projection
$\pi_i$ of a product stratum $\Delta_\Gamma\cong\prod_j \M_{g_j,n_j}$, then the
complementary product $\prod_{j\neq i}\M_{g_j,n_j}$ contains a complete curve.
Consequently:
\begin{enumerate}
\item[(a)] if each factor $\M_{g_j,n_j}$ with $j\ne i$ contains no complete curves,
the hypotheses of Claim~\ref{claim:P} are contradictory in the closed algebraic
category (vacuity of the claim for such strata);
\item[(b)] a smooth curve minus a nonempty finite set is affine (classical), so
$\M_{g',1}\to\M_{g'}$ has fibers containing no positive-dimensional complete
subvarieties; by the bound of Diaz \cite{Diaz} (complete subvarieties of $\M_{g'}$
have dimension $\le g'-2$), the spaces $\M_{g'}$ for $g'\le 2$ and $\M_{g',1}$ for
$g'\le2$ contain no complete curves; $\M_{0,n}$ embeds into an affine space by the
cross-ratio maps, and $\M_{1,1}\cong \C^j$ is affine; hence (a) applies whenever the
complementary factor has genus $\le 2$ (e.g.\ $g=3$ with a one-separating-edge
stratum, where the complementary genus is $1$ or $2$).
\end{enumerate}
\end{enumerate}
\end{proposition}

\begin{proof}
(1): $\overline{V}$ is a compact analytic subset of the projective variety
$\overline{\M}_g$, hence algebraic. (2): $\partial V$ is compact iff it is closed in
$\overline{V}$; if not, its limit points in $\overline{V}$ lie in
$\overline{\Delta_\Gamma}\setminus\Delta_\Gamma$, i.e.\ in deeper strata. When
$\partial V$ is compact, it is a compact subset of the Hausdorff space
$\overline{\M}_g$, locally analytic, hence a closed analytic subset of
$\overline{\M}_g$, and therefore algebraic and complete. (3): A complete curve
$\partial V$ contained in the fiber $\pi_i^{-1}(Y)\cong\prod_{j\ne i}\M_{g_j,n_j}$
projects to complete subvarieties in each factor; since $\dim\partial V=1$, at least
one projection is one-dimensional (otherwise $\partial V$ would be a point), so that
factor contains a complete curve. For (b): if $K\subset \M_{g',1}$ ($g'\le 2$) is
complete, its image in $\M_{g'}$ is complete, hence a point by Diaz's bound (and,
for $g'=0,1$, by the affine embeddings recalled above); then $K$ lies in a fiber, a
punctured curve, which is affine and therefore contains no positive-dimensional
complete subvarieties; so $K$ is a point. The remaining assertions are immediate.
\end{proof}

\begin{corollary}\label{cor:burden}
Any correct proof of Claim~\ref{claim:P}, in any of the admissible readings, must
derive the constancy of the holomorphic map $Y$ solely from the hypothesis that the
central fiber $V_0$ is totally geodesic for the Teichm\"uller metric.
\end{corollary}

\begin{proof}
Combine Lemma~\ref{lem:hol} and Corollary~\ref{cor:equiv} with
Theorem~\ref{thm:model} and Remark~\ref{rem:scope}: the remaining hypotheses hold in
configurations with $Y$ non-constant, so they cannot imply constancy.
\end{proof}

\section{First-order analysis: there is no infinitesimal rigidity}\label{sec:inflect}

We now show that, at the infinitesimal level, nothing obstructs the motion of the
collapse point. We work in the following abstract setting, which models a neighborhood
of $\overline{V_0}$ in a projective model of $\overline{\M}_g$.

\begin{setup}\label{setup:abstract}
Let $X$ be a smooth complex projective variety, $D\subset X$ a smooth closed
subvariety, $B$ a complex manifold, and $\pi:D\to B$ a holomorphic map which is the
projection of a product $D\cong B\times F$ over a neighborhood of a point $b_0\in B$.
Let $S\subset X$ be a smooth closed subvariety of dimension $2$ such that $S$
intersects $D$ transversally along a smooth (possibly disconnected) compact curve
$C=S\cap D$, and suppose $\pi(C)=\{b_0\}$. Let $N:=N_{S/X}$ denote the normal bundle.
\end{setup}

\begin{lemma}\label{lem:iso}
Transversality induces a canonical isomorphism $\alpha:N|_C\stackrel{\sim}{\longrightarrow}
N_{C/D}=(T_D|_C)/T_C$. Consequently, restriction gives a map
$r:H^0(S,N)\to H^0(C,N_{C/D})$.
\end{lemma}

\begin{proof}
Along $C$ we have $T_X=T_S+T_D$ (transversality) and $T_S\cap T_D=T_C$. Hence
$N|_C=T_X|_C/T_S|_C=(T_S+T_D)/T_S\cong T_D/(T_S\cap T_D)=T_D/T_C=N_{C/D}$, naturally
in the local coordinates.
\end{proof}

\begin{lemma}\label{lem:motion}
Let $(S_\epsilon)_{|\epsilon|<\epsilon_0}$ be a holomorphic arc of submanifolds of $X$
with $S_0=S$, transverse to $D$, with first-order normal field
$v\in H^0(S,N)$; i.e., the arc's tangent in the Hilbert scheme $\Hilb(X)$ at $[S]$
is $v$. Then the trace family $C_\epsilon:=S_\epsilon\cap D$ is a holomorphic family
of curves in $D$ whose first-order normal field in $N_{C/D}$ equals
$r(v)=\alpha(v|_C)$.
\end{lemma}

\begin{proof}
Transversality is an open condition, so $C_\epsilon$ is well defined and varies
holomorphically (the intersection of a flat family with a fixed smooth subvariety,
locally cut out by a single non-zero-divisor equation $u$, remains flat;
equivalently, work in local coordinates). Choose coordinates
$(c,u,\mathbf z)$ near a point of $C$ with $D=\{u=0\}$, $S=\{\mathbf z=0\}$,
$C=\{u=0,\ \mathbf z=0\}$. The arc has the local form
$S_\epsilon=\{\mathbf z=\epsilon\,\vec v(c,u)+O(\epsilon^2)\}$, where
$\vec v\,\partial_{\mathbf z}$ represents $v$. Intersecting with $\{u=0\}$ gives
$C_\epsilon=\{u=0,\ \mathbf z=\epsilon\,\vec v(c,0)+O(\epsilon^2)\}$, whose
first-order normal field in $D$ is $\vec v(c,0)\partial_{\mathbf z}$; this is exactly
$\alpha(v|_C)$. Globally, these local computations glue.
\end{proof}

\begin{theorem}[Infinitesimal flexibility]\label{thm:inflect}
In Setup~\ref{setup:abstract}, assume
\[
H^1(S,N)=0,\qquad H^1(S,N\otimes_{\mathcal O_S}\mathcal O_S(-C))=0 .
\]
Then for every $\xi\in T_{b_0}B\setminus\{0\}$ there exists a holomorphic arc
$(S_\epsilon)$ of submanifolds of $X$ with $S_0=S$, tangent velocity $v=v(\xi)\in
H^0(S,N)$, transverse to $D$, such that the first-order trace deformation
$w:=r(v)\in H^0(C,N_{C/D})$ satisfies
\[
d\pi(w)\equiv \xi\quad\text{as a constant section of } T_{b_0}B\otimes \mathcal O_C.
\]
Consequently, if the traces of an arc realizing $v(\xi)$ remain collapsed
(i.e.\ $\pi(C_\epsilon)=\{Y_\epsilon\}$ is a single point for all small
$\epsilon\ne0$), then $Y'(0)=\xi\ne 0$.
\end{theorem}

\begin{proof}
\emph{Step 1 (restriction is surjective).} On the smooth surface $S$, the curve $C$
is an effective Cartier divisor, so the standard sequence
$0\to N(-C)\to N\to N|_C\to0$ is exact, and the hypothesis
$H^1(S,N(-C))=0$ makes $H^0(S,N)\to H^0(C,N|_C)$ surjective. Composing with the
isomorphism $\alpha$ of Lemma~\ref{lem:iso}, the map
$r:H^0(S,N)\to H^0(C,N_{C/D})$ is surjective.

\emph{Step 2 (the desired trace velocity exists).} Choose a holomorphic vector field
$\xi_B$ on a neighborhood of $b_0$ in $B$ with $\xi_B(b_0)=\xi$. Since $\pi$ is the
projection of a product near the (compact) $C$, the pullback $\tilde\xi$ of $\xi_B$
is a holomorphic vector field on a neighborhood of $C$ in $D$, tangent to $D$, with
$d\pi(\tilde\xi)=\xi_B\circ\pi$. Let $w_\xi\in H^0(C,N_{C/D})$ be the class of
$\tilde\xi|_C$. Because $\pi|_C\equiv b_0$, we have $d\pi(T_C)=0$, so $d\pi$
descends to a map $N_{C/D}\to T_{b_0}B\otimes\mathcal O_C$, and
$d\pi(w_\xi)=\xi_B(b_0)\equiv\xi$, a constant section.

\emph{Step 3 (lifting and integrating).} Choose $v\in H^0(S,N)$ with $r(v)=w_\xi$.
Since $S\subset X$ is a local complete intersection and $H^1(S,N)=0$, the Hilbert
scheme $\Hilb(X)$ is smooth at $[S]$ with tangent space $H^0(S,N)$ and obstruction
space $H^1(S,N)$ \cite{Sernesi, HartshorneDef, Nitsure}. At a smooth point every
tangent vector is the velocity of a holomorphic arc; choose an arc $(S_\epsilon)$
with velocity $v$, and shrink so that all $S_\epsilon$ are transverse to $D$.

\emph{Step 4 (conclusion).} By Lemma~\ref{lem:motion}, the trace family
$C_\epsilon=S_\epsilon\cap D$ has first-order normal field $w_\xi$, so
$d\pi(\dot C)\equiv\xi$. If $\pi(C_\epsilon)$ is a single point $Y_\epsilon$ for all
small $\epsilon$, differentiating the composite
$\epsilon\mapsto \pi\circ(\text{trace})$ at $\epsilon=0$ gives
$Y'(0)=\xi\ne 0$.
\end{proof}

\begin{proposition}[The rigidity that Claim $\mathsf{P}$ would have to establish]
\label{prop:R}
Let $\mathfrak C$ denote the class of all tuples satisfying the hypotheses of
Claim~\ref{claim:P} (in the Standing Interpretation). Assume $\mathfrak C\ne\emptyset$,
and let $(g,\Gamma,i,V_0,(V_t))\in\mathfrak C$ be such that $S:=\overline{V_0}$ is a
smooth compact surface meeting $D:=\Delta_\Gamma$ transversely along
$C=S\cap D$ and satisfying the two $H^1$-vanishing hypotheses of
Theorem~\textup{\ref{thm:inflect}} (with $X$ a smooth projective model of
$\overline{\M}_g$ along $S$, $B=\M_{g_i,n_i}$, $\pi=\pi_i$). Then Claim
$\mathsf{P}$ \emph{implies} the following assertion \textup{(R)}:
\begin{quote}
\textup{(R)}\; For every $\xi\in T_{Y_0}\M_{g_i,n_i}\setminus\{0\}$ and every arc
$(S_\epsilon)$ as in Theorem~\textup{\ref{thm:inflect}} with velocity $v(\xi)$ and
with traces collapsed (and with the incidence conditions (i), (ii) propagating, an
open condition on arcs near the constant arc), the collapse point is constant:
$Y_\epsilon\equiv Y_0$.
\end{quote}
In particular, Claim $\mathsf{P}$ entails that \emph{none} of the first-order
deformations produced by Theorem~\textup{\ref{thm:inflect}} integrates inside the
class of collapsed traces.
\end{proposition}

\begin{proof}
An arc $(S_\epsilon)$ of surfaces, transverse to the boundary and with collapsed
traces, yields a holomorphic family of submanifolds
$V_\epsilon:=S_\epsilon\cap \M_g$ satisfying the problem's footnote and, for arcs
near the constant arc, hypotheses (i) and (ii): proper holomorphic submersion trace
families are smoothly trivial by Fischer--Grauert \cite{FG} (hypothesis (i)), and the
incidence pattern (ii) is preserved under small perturbations whenever the
intersections with the deeper strata are transverse. For such a family,
Claim~\ref{claim:P} asserts $Y_\epsilon\equiv Y_0$. Theorem~\ref{thm:inflect}
exhibits, for each $\xi\ne0$, the required tangent data, so (R) holds.
\end{proof}

\begin{remark}\label{rem:nomech}
Assertion (R) is an \emph{extrinsic, deformation-theoretic} rigidity: it demands that
certain first-order deformations---which exist unconditionally at the level of
$H^0(S,N)$ under the stated (generic, open) hypotheses---fail to integrate within a
certain closed locus of the Hilbert scheme. None of the hypotheses of
Claim~\ref{claim:P} other than total geodesy bears on integrability, and by
Theorem~\ref{thm:model} those hypotheses are satisfied in situations where $Y$ moves.
Thus the whole burden falls on the total geodesy of the single central fiber. No
known theorem---classical or from the modern theory of Teichm\"uller dynamics---converts
an intrinsic metric condition on one member of a family into the failure of
integrability of prescribed first-order deformations of the ambient subvariety. See
Section~\ref{sec:knowledge}.
\end{remark}

\section{The total geodesy hypothesis: state of knowledge}\label{sec:knowledge}

\subsection{What is known about totally geodesic subvarieties of $\M_g$}

We collect the relevant facts, with references.

\begin{enumerate}
\item The Teichm\"uller space $(\T_g,d_T)$ with the Teichm\"uller metric is a complete
geodesic metric space; distances are realized by extremal Teichm\"uller maps, and
geodesic segments are governed by holomorphic quadratic differentials (Teichm\"uller's
theorem; see \cite{Hubbard06, IT}). Royden proved that the Teichm\"uller metric
coincides with the Kobayashi metric of $\T_g$ \cite{Royden}.
\item An algebraic curve $C\subset \M_g$ is totally geodesic for $d_T$ if and only if
it is a \emph{Teichm\"ller curve}, i.e.\ the algebraic image of a Teichm\"uller disc
whose stabilizer (Veech group) is a lattice; this is the standard characterization of
Teichm\"uller curves (see \cite{McMullen, Moeller, EMM}).
\item Teichm\"uller curves exist in abundance (e.g.\ square-tiled curves), and a great
deal is known about their arithmetic (real multiplication, fields of definition,
finiteness phenomena: \cite{Moeller, EMM, McMullen}).
\item \emph{In dimensions $\ge2$, however, no totally geodesic complex subvariety of
any $\M_g$ is known.} To the best of the author's knowledge, every totally geodesic
subvariety of $\M_g$ with respect to the Teichm\"uller metric that has ever been
constructed has dimension $\le 1$, and the existence of a totally geodesic complex
\emph{surface} in some $\M_g$ appears to be an open problem. The classification
program of affine invariant submanifolds and orbit closures
\cite{EMM, Moeller} singles out Teichm\"uller curves as the only known ``algebraic''
projections, and no higher-dimensional totally geodesic candidate has been produced
in the decades since.
\end{enumerate}

\begin{remark}
The hypotheses of Claim~\ref{claim:P} require $\dim V_0=2$ exactly: for a
one-dimensional totally geodesic $V_0$ the traces $\overline{V_t}\cap\Delta_\Gamma$
would be zero-dimensional and the hypothesis $\dim\partial V_t=1$ could not hold.
Thus the objects to which Claim~\ref{claim:P} applies are precisely of the unknown
kind described in (4).
\end{remark}

\subsection{No transfer from intrinsic geodesy to extrinsic deformation theory}

Total geodesy of $V_0$ is a statement comparing the intrinsic geodesics of (the
Teichm\"uller lift of) $V_0$ with those of $\M_g$: it constrains the metric
geometry of a single fiber. The deformation theory of $V_0$ as a submanifold,
by contrast, is governed by the extrinsic data
$H^0(\overline{V_0},N_{\overline{V_0}/\overline{\M}_g})$
(\cite{Sernesi, HartshorneDef}); the collapse map $Y$ is holomorphic
(Lemma~\ref{lem:hol}) and its infinitesimal behavior is dictated by the restriction
maps analyzed in Section~\ref{sec:inflect}. These two regimes are completely
disjoint:

\begin{itemize}
\item In dimension one, where the theory of totally geodesic subvarieties (Teichm\"uller
curves) is essentially complete, no analogous extrinsic rigidity is known: nothing in
the theory prevents a Teichm\"uller curve from possessing extrinsic deformations whose
boundary traces move; whether such deformations exist for a given curve is a question
about its normal bundle, not about total geodesy.
\item The hypotheses (i), (ii) of Claim~\ref{claim:P} are pure regularity conditions
and, by Theorem~\ref{thm:model}, are satisfied in configurations where $Y$ moves.
\item No result in the literature---classical deformation theory, the theory of
Teichm\"uller curves, or the Eskin--Mirzakhani--Mohammadi orbit-closure
machinery---implies anything like assertion (R) of Proposition~\ref{prop:R}.
\end{itemize}

\subsection{The dichotomy}

\begin{theorem}[Dichotomy]\label{thm:dichotomy}
Let $\mathfrak C$ be the class of tuples satisfying the hypotheses of
Claim~\ref{claim:P}. Then:
\begin{enumerate}
\item If $\mathfrak C=\emptyset$, then Claim~\ref{claim:P} holds vacuously; however,
establishing $\mathfrak C=\emptyset$ includes in particular proving that no totally
geodesic complex surface exists in any $\M_g$, which is an open problem
(Section~\textup{\ref{sec:knowledge}}).
\item If $\mathfrak C\ne\emptyset$, then Claim~\ref{claim:P} holds only if, for every
admissible totally geodesic surface $V_0$ satisfying the mild genericity hypotheses of
Theorem~\textup{\ref{thm:inflect}}, the rigidity assertion \textup{(R)} of
Proposition~\textup{\ref{prop:R}} holds: all of the first-order deformations with
non-zero infinitesimal collapse velocity fail to integrate within the class of
collapsed traces. This is a rigidity phenomenon of a type not implied by any known
theorem, and it is false in the absence of the total-geodesy hypothesis
(Theorem~\textup{\ref{thm:model}}).
\end{enumerate}
In particular, in neither branch of the dichotomy can a complete rigorous proof of
Claim~\ref{claim:P} be given at present, and no such proof can be obtained from the
stated hypotheses without first resolving an open problem or establishing a major new
rigidity theorem.
\end{theorem}

\begin{proof}
(1) is immediate from the definitions and Section~\ref{sec:knowledge}(4). For (2),
let $(g,\Gamma,i,V_0,(V_t))\in\mathfrak C$. If $V_0$ (more precisely
$S=\overline{V_0}$) satisfies the genericity hypotheses of
Theorem~\ref{thm:inflect}, then by Proposition~\ref{prop:R} the truth of
Claim~\ref{claim:P} for this tuple forces (R). If no member of $\mathfrak C$ satisfies
those hypotheses, then Claim~\ref{claim:P} is true only if, a fortiori, every totally
geodesic surface with the stated boundary structure fails the genericity hypotheses
or otherwise obstructs all the relevant arcs---again a nontrivial, unprecedented
statement about the unknown objects of Section~\ref{sec:knowledge}(4). The final
assertion follows from Corollary~\ref{cor:burden}.
\end{proof}

\section{Conclusion}\label{sec:conc}

\begin{theorem}[No-go]\label{thm:nogo}
The assertion of Claim~\textup{\ref{claim:P}} is not a theorem derivable from the
hypotheses stated in the problem, and at present no complete rigorous proof of it
can be given. More precisely:
\begin{enumerate}
\item the conclusion is equivalent to the constancy of a canonically defined
holomorphic map $Y:\Delta\to\M_{g_i,n_i}$ (Lemma~\textup{\ref{lem:hol}});
\item all hypotheses except the total geodesy of $V_0$ are compatible with
non-constant $Y$ (Theorem~\textup{\ref{thm:model}});
\item under mild genericity hypotheses, first-order deformations with non-zero
infinitesimal collapse velocity always exist (Theorem~\textup{\ref{thm:inflect}}),
so the claim would have to be established through the unprecedented rigidity
\textup{(R)} of Proposition~\textup{\ref{prop:R}};
\item the objects to which the claim applies are not known to exist: the existence
of a totally geodesic complex surface in any $\M_g$ is an open problem
(Section~\textup{\ref{sec:knowledge}});
\item hence, by Theorem~\textup{\ref{thm:dichotomy}}, the claim is either vacuous
with open vacuity, or it asserts an unknown and hitherto undemonstrated rigidity for
hypothetical objects.
\end{enumerate}
Consequently, any purported complete proof of Claim~\textup{\ref{claim:P}} from the
stated hypotheses must contain an error, or must silently resolve one of the open
problems identified here.
\end{theorem}

\begin{remark}[What might be true instead]
We record, for completeness, the variants under which a statement of this shape could
become a well-posed problem:
\begin{itemize}
\item If one assumes that \emph{all} fibers $V_t$ are totally geodesic, the question
becomes a rigidity question within Teichm\"uller dynamics; even then, no theorem of
the asserted form is known, and the problem explicitly forgoes this hypothesis.
\item If one strengthens the collapse condition by arithmetic requirements (e.g.\
that the collapsed point $Y_t$ carry a fixed field of definition or a fixed
isomorphism type constrained by real multiplication), the regularity hypotheses (i),
(ii) still cannot produce constancy; such hypotheses are absent from the problem.
\item Without the total-geodesy hypothesis the asserted statement is simply false
in the category of Theorem~\ref{thm:model}.
\end{itemize}
\end{remark}

\begin{remark}[On the requested task]
We were asked to ``work very hard and not return until a complete proof is
achieved.'' We have instead worked hard to determine whether a complete proof can
exist. The outcome of that analysis, given with full proofs of all supporting
statements above, is negative: the assertion in the question is either vacuous by
virtue of an open problem or expresses an unknown rigidity; in neither case can the
requested proof honestly be supplied, and we decline to fabricate one.
\end{remark}

\begin{thebibliography}{99}\small

\bibitem{DM} P. Deligne and D. Mumford, \emph{The irreducibility of the space of
curves of given genus}, Inst. Hautes \'Etudes Sci. Publ. Math. \textbf{36} (1969),
75--109.

\bibitem{Knudsen} F. F. Knudsen, \emph{The projectivity of the moduli space of stable
curves, III}, Math. Scand. \textbf{52} (1983).

\bibitem{ACGII} E. Arbarello, M. Cornalba, and P. A. Griffiths (with J. D. Harris),
\emph{Geometry of Algebraic Curves, Volume II}, Grundlehren der math.
Wissenschaften \textbf{268}, Springer, 2011.

\bibitem{Hubbard76} J. Hubbard, \emph{Sur les sections analytiques de la courbe
universelle de Teichm\"uller}, Amer. J. Math. \textbf{98} (1976), 327--364.

\bibitem{Hubbard06} J. Hubbard, \emph{Teichm\"uller Theory and Applications to
Geometry, Topology, and Dynamics, Vol.~1}, Matrix Editions, 2006.

\bibitem{IT} Y. Imayoshi and M. Taniguchi, \emph{An Introduction to Teichm\"uller
Spaces}, Springer, 1992.

\bibitem{Royden} H. L. Royden, \emph{Automorphisms and isometries of Teichm\"uller
space}, in: Advances in the Theory of Riemann Surfaces, Ann. of Math. Studies
\textbf{66}, Princeton Univ. Press, 1971.

\bibitem{McMullen} C. T. McMullen, \emph{Teichm\"uller curves in genus two: the
decagon and the Bolza curve}, J. Amer. Math. Soc. \textbf{18} (2005).

\bibitem{Moeller} M. M\"oller, \emph{Variations of Hodge structures of Teichm\"uller
curves}, J. Amer. Math. Soc. \textbf{19} (2006), 327--344.

\bibitem{EMM} A. Eskin, M. Mirzakhani, and A. Mohammadi, \emph{Isolation,
equidistribution, and orbit closures for the $SL(2,\mathbb R)$ action on moduli
space}, Ann. Sci. \'Ec. Norm. Sup\'er. \textbf{47} (2014).

\bibitem{FG} W. Fischer and H. Grauert, \emph{Locally trivial deformations of complex
structures}, Math. Ann. \textbf{160} (1965), 72--88.

\bibitem{Sernesi} E. Sernesi, \emph{Deformations of Algebraic Schemes}, Grundlehren
der math. Wissenschaften \textbf{336}, Springer, 2006.

\bibitem{HartshorneDef} R. Hartshorne, \emph{Deformation Theory}, Graduate Texts in
Mathematics \textbf{257}, Springer, 2010.

\bibitem{Nitsure} N. Nitsure, \emph{Construction of Hilbert and Quot schemes}, in:
Fundamental Algebraic Geometry: FGA Explained, Math. Surveys and Monographs
\textbf{123}, Amer. Math. Soc., 2005.

\bibitem{Chirka} E. M. Chirka, \emph{Complex Analytic Sets}, Kluwer, 1989.

\bibitem{Demailly} J.-P. Demailly, \emph{Complex Analytic and Differential Geometry},
book draft, available online.

\bibitem{Diaz} S. Diaz, \emph{A bound on the dimensions of complete subvarieties of
$\mathcal M_g$}, Duke Math. J. \textbf{51} (1984).

\bibitem{Kodaira} K. Kodaira, \emph{A certain type of irregular algebraic surfaces},
J. Analyse Math. \textbf{19} (1967), 207--215.

\end{thebibliography}

\end{document}
```

\medskip

\noindent\textbf{Summary of the answer.}
The document above does not contain the requested proof, because a complete rigorous
proof of the asserted statement cannot be given. What it contains instead, with full
proofs, is the following analysis. (1)~Under the natural reading of the problem, the
collapse point is a well-defined holomorphic function of $t$ (Fischer--Grauert
descent), so the claim is a genuine rigidity statement. (2)~An explicit local model
satisfies every hypothesis except total geodesy of the central fiber while the
collapse point moves freely; hence all the regularity hypotheses (i), (ii) are
powerless, and any proof would have to derive the constancy from total geodesy alone.
(3)~First-order deformation theory (normal-bundle sequence plus Hilbert-scheme
integration) shows that, under mild genericity, deformations with non-zero
infinitesimal velocity of the collapse point always exist, so the claim would require
an unprecedented ``failure of integrability'' rigidity. (4)~No totally geodesic
complex surface in any $\mathcal M_g$ has ever been constructed, so the objects to
which the claim applies are not known to exist. The claim is therefore either
vacuous (with vacuity equivalent to an open problem) or it asserts an unknown
rigidity; in neither branch does a proof from the stated hypotheses exist, and any
purported proof must contain an error.